% ===========================================================================
% Vue de déploiement — dérivée de docker-compose.yml et frontend/nginx.conf,
% relevés le 2026-09-22.
%
% Trois faits que ce diagramme doit rendre INCONTOURNABLES :
%   1. le conteneur backend ne publie AUCUN port vers l'hôte — il n'est
%      joignable qu'à travers nginx, depuis le réseau interne ;
%   2. il n'y a AUCUNE terminaison TLS : le point d'entrée est en clair ;
%   3. il n'existe AUCUN environnement cible — cette composition est celle du
%      poste de développement.
%
% Largeur tenue sous 150 mm : au-delà, la figure déborde de la justification.
% ===========================================================================
\begin{tikzpicture}[node distance=7mm]

  % ─── Hors frontière ─────────────────────────────────────────────────────
  \node[pqactor, text width=24mm] (nav) {Navigateur\\{\scriptsize poste client}};
  \node[pqcaptionwarn, above=2mm of nav] {hors frontière};

  % ─── Conteneurs applicatifs ─────────────────────────────────────────────
  \node[pqbox, right=11mm of nav, text width=35mm] (fe)
    {\textbf{piqbit-frontend}\\[2pt]
     {\scriptsize nginx 1.27, port \texttt{80}}\\
     {\scriptsize bundle Angular statique}\\
     {\scriptsize relais \texttt{/api/} et WebSocket}};

  \node[pqcore, right=11mm of fe, text width=36mm] (be)
    {\textbf{piqbit-backend}\\[2pt]
     {\scriptsize\mdseries FastAPI/uvicorn, port \texttt{8000}}\\
     {\scriptsize\mdseries utilisateur non privilégié}\\
     {\scriptsize\mdseries migration au démarrage}};

  % ─── Services d'infrastructure ──────────────────────────────────────────
  \node[pqstore, below=15mm of fe, xshift=6mm, minimum width=25mm] (pg)
    {\textbf{postgres}\\pgvector pg16\\\texttt{5434:5432}};
  \node[pqstore, right=8mm of pg, minimum width=25mm] (rd)
    {\textbf{redis}\\7-alpine\\\texttt{6379:6379}};
  \node[pqstore, right=8mm of rd, minimum width=25mm] (vol)
    {volume\\\texttt{uploads}\\{\scriptsize CV déposés}};

  % ─── Frontière du réseau Docker ─────────────────────────────────────────
  \begin{scope}[on background layer]
    \node[pqboundary, fit=(fe)(be)(pg)(rd)(vol), inner sep=9pt] (net) {};
  \end{scope}
  \node[pqcaption, anchor=south west, xshift=2mm] at (net.north west)
    {Réseau Docker interne — profil \texttt{app}};

  % ─── Liens ──────────────────────────────────────────────────────────────
  \draw[pqflow] (nav.east) -- node[pqlbl, align=center, pos=0.5]
    {HTTP\\\texttt{:8080}} (fe.west);
  \draw[pqflow] (fe.east) -- node[pqlbl, align=center, pos=0.5]
    {\texttt{proxy\_pass}} (be.west);

  % Départ depuis trois points distincts du bord inférieur : un départ unique
  % produisait trois traits superposés sur la première moitié du parcours.
  \draw[pqflow] ([xshift=-12mm]be.south) -- ([xshift=-12mm]be.south |- pg.north) -- (pg.north);
  \draw[pqflow] (be.south)               -- (be.south |- rd.north)               -- (rd.north);
  \draw[pqflow] ([xshift=12mm]be.south)  -- ([xshift=12mm]be.south |- vol.north) -- (vol.north);

  % ─── Constats à déclarer sans détour ────────────────────────────────────
  \node[pqwarn, below=8mm of net.south, anchor=north, text width=132mm]
    {\textbf{Trois constats à déclarer, pas à contourner.}
     \emph{Aucune terminaison TLS} dans cette composition : le point d'entrée
     \texttt{:8080} est en clair. Le conteneur backend \emph{ne publie aucun
     port} vers l'hôte — il n'est atteignable qu'à travers nginx, ce qui est
     une propriété voulue. \emph{Aucun environnement d'hébergement cible}
     n'existe à ce jour, et aucune procédure de retour arrière n'est définie :
     la composition ci-dessus est celle du poste de développement.};

\end{tikzpicture}
